\begin{examplebox}
\textbf{例 1\quad 数值求最大顶角（原例题纠错）。}

椭圆 $\frac{x^2}{4}+y^2=1$，$a=2,b=1,c^2=3$。
求 $\alpha=\angle A_1PA_2$ 的最大值。

\textbf{解：}在 $P=(0,\pm1)$ 取到最大值；
\[
\boxed{\alpha_{\max}=2\arctan2
\approx126.87^\circ},\qquad
\boxed{\tan\alpha_{\max}=-\frac43}.
\]
原稿的 $\arctan(4/3)\approx53.13^\circ$ 是\textbf{补角对应的锐角}，
不是这里的三角形顶角。

\medskip
\textbf{例 2\quad 角度存在性：为什么 $60^\circ$ 永远不可能？}

问：椭圆上是否存在非端点 $P$，使
$\angle A_1PA_2=60^\circ$ 或 $120^\circ$？

\textbf{解：}由于 $\alpha>90^\circ$，$60^\circ$ 在任何非圆椭圆上
\textbf{都不可能}。若要求 $120^\circ$ 存在，等价于
\[
2\arctan\frac ab\ge120^\circ
\iff \frac ab\ge\sqrt3
\iff \frac ba\le\frac1{\sqrt3}.
\]
由 $e=c/a$、$b^2/a^2=1-e^2$，得
\[
\boxed{\sqrt{\frac23}\le e<1}.
\]
等号时只能在短轴端点取得 $120^\circ$；
严格大于时则有四个对称的解点。

\medskip
\textbf{例 3\quad 已知顶角 $120^\circ$，反求 $P$ 的全部坐标。}

仍取椭圆 $\frac{x^2}{4}+y^2=1$，已知
$\alpha=120^\circ$，求所有可能的 $P$。

\textbf{解：}此时 $\delta=\alpha-90^\circ=30^\circ$，
利用 $\tan\delta=3|y|/4$，
\[
|y|=\frac4{3\sqrt3}=\frac{4\sqrt3}{9},
\qquad
|x|=2\sqrt{1-y^2}
=2\sqrt{1-\frac{16}{27}}
=\frac{2\sqrt{33}}9.
\]
所以\textbf{共有四点}：
\[
\boxed{P=\biggl(\pm\frac{2\sqrt{33}}9,\,
\pm\frac{4\sqrt3}9\biggr)}
\quad\text{（两处 $\pm$ 独立取号）。}
\]
检查 $|y|<1$，符合存在性条件。

\medskip
\textbf{例 4\quad 给出斜率积，反求最大顶角。}

已知 $k_{PA_1}k_{PA_2}=-\frac14$，
其中 $P$ 在长轴水平的非圆椭圆上。
由斜率积定值 $-b^2/a^2=-1/4$，得 $a/b=2$，
于是
\[
\boxed{\alpha_{\max}=2\arctan2\approx126.87^\circ},
\quad\boxed{e=\sqrt3/2}.
\]
\textbf{辨析：}斜率积恒定，\textbf{不代表夹角恒定}；
两个斜率各自仍随点 $P$ 改变。

\medskip
\textbf{例 5\quad 面积与角度最值：两题为何同解？}

设 $S$ 为 $\triangle A_1PA_2$ 的面积。
因为 $S=a|y|$，$0<|y|\le b$，故
\[
\boxed{0<S\le ab},\quad
S_{\max}=ab\text{ 在 }P=(0,\pm b)\text{ 取得}.
\]
同时
\[
\tan(\alpha-90^\circ)=\frac{c^2}{2a^2b^2}S.
\]
因此求 $S$ 最大与求 $\alpha$ 最大，\textbf{本质都是求 $|y|$ 最大}。
面积无最小值，下确界为零，也不能说长轴端点处取到最小面积
（因为 $P$ 被题设排除）。

\medskip
\textbf{例 6\quad 已知最大顶角，反求离心率。}

若 $\alpha_{\max}=120^\circ$，则
$\arctan(a/b)=60^\circ$，$a/b=\sqrt3$，所以
\[
\boxed{e=\sqrt{1-\frac{b^2}{a^2}}
=\sqrt{\frac23}}.
\]
\textbf{迁移：}由最大顶角可反求轴比，
进而反求离心率；这属于“角度量 $\to$ 形状参数”。

\medskip
\textbf{例 7\quad 同一椭圆上的等角点个数。}

设已知固定椭圆的 $\alpha_{\max}=130^\circ$，
研究目标角 $\alpha_0$ 的点数：
\begin{itemize}
\item $\alpha_0=100^\circ$：四个对称点；
\item $\alpha_0=130^\circ$：两个短轴端点；
\item $\alpha_0=90^\circ$：零个点（只能趋近，不能取到）；
\item $\alpha_0=135^\circ$：零个点（超过最大值）。
\end{itemize}
\textbf{识别信号：}先判断 $90^\circ<\alpha_0\le\alpha_{\max}$，
再决定是否需要反求坐标。

\medskip
\textbf{例 8\quad 从 C003 焦角直接联想到 C004 顶角。}

椭圆 $a=2,b=1,c=\sqrt3$；
若某点 $P$ 的焦角为 $\theta$，则
\[
\tan(\alpha-90^\circ)
=\frac{\sqrt3}{4}\tan\frac\theta2.
\]
特别在 $P=(0,\pm1)$，焦角半角正切为 $c/b=\sqrt3$，
故 $\tan(\alpha-90^\circ)=3/4$，得到
$\alpha=90^\circ+\arctan(3/4)\approx126.87^\circ$。
\textbf{提醒：}这里 $\theta$ 与 $\alpha$ 是两个不同的角。
\end{examplebox}
